\begin{trapbox}
\textbf{误区一：分母较大就是 $a^2$。} 这是椭圆的读法；双曲线应看正项。
\begin{center}\begin{tikzpicture}[font=\small,>=stealth]
\draw[black!30,rounded corners=2pt] (0,0) rectangle (14.4,8.2);
\node[anchor=west,font=\small\bfseries] at (.5,7.72) {实半轴 $a$ 由正项决定；分母较大也可能是 $b^2$};
\begin{scope}[shift={(3.6,4)},scale=.34]
\draw[black!35,->] (-7,0)--(7,0) node[right] {$x$};\draw[black!35,->] (0,-7)--(0,7) node[above] {$y$};
\draw[blue,very thick,domain=-5.5:5.5,samples=101,variable=\t] plot ({-1*3*sqrt(1+\t*\t/16)},\t);
\fill (-5,0) circle (2pt) node[below=4pt] {$F_{1}$};
\draw[blue,very thick,domain=-5.5:5.5,samples=101,variable=\t] plot ({1*3*sqrt(1+\t*\t/16)},\t);
\fill (5,0) circle (2pt) node[below=4pt] {$F_{2}$};
\end{scope}
\node at (3.6,1.1) {$x^2/9-y^2/16=1$};
\begin{scope}[shift={(10.8,4)},scale=.34]
\draw[black!35,->] (-7,0)--(7,0) node[right] {$x$};\draw[black!35,->] (0,-7)--(0,7) node[above] {$y$};
\draw[blue,very thick,domain=-5.5:5.5,samples=101,variable=\t] plot (\t,{-1*3*sqrt(1+\t*\t/16)});
\fill (0,-5) circle (2pt) node[right=5pt] {$F_{1}$};
\draw[blue,very thick,domain=-5.5:5.5,samples=101,variable=\t] plot (\t,{1*3*sqrt(1+\t*\t/16)});
\fill (0,5) circle (2pt) node[right=5pt] {$F_{2}$};
\end{scope}
\node at (10.8,1.1) {$y^2/9-x^2/16=1$};

\node[anchor=west] at (.7,.55) {同尺度；两图都是 $a=3,b=4,c=5$；焦点轴随正项交换};
\end{tikzpicture}\end{center}

两图的较大分母都是 $16=b^2$；横向与纵向只交换坐标，不交换实、虚半轴的含义。

\medskip
\textbf{误区二：沿用椭圆的减号。} 椭圆为 $a^2=b^2+c^2$，双曲线为 $c^2=a^2+b^2$。
同焦点图已显示：椭圆实顶点在焦点外侧，双曲线实顶点在焦点内侧。
因此双曲线算出的 $c$ 若不大于 $a,b$，应回查。

\medskip
\textbf{误区三：把虚半轴端点或矩形角点当作曲线点。}
横向双曲线与 $y$ 轴没有交点，$(a,b)$ 代入左边为 $0$，不是 $1$。
它们只用于辅助矩形和渐近线构造；参数图的斜边是 $OQ$。

\medskip
\textbf{误区四：焦距与半焦距混用，或开方写正负。}
$c>0$ 是长度，焦距才是 $2c$；正负号用于焦点坐标。
已知 $a,c$，若 $c\le a$，则 $b^2\le0$，不满足本条非退化双曲线条件。
\end{trapbox}
